\chapter{Implementation}
\label{cha:implementation}

\epigraph{Ideas are easy, Implementation is hard.}{Guy Kawasaki}

Most of the implementation details and general decisions made for developing a
real-world functional project, based on the design architecture presented in chapter
\ref{cha:architecture}, are discussed in this chapter. It is organized in such a
way that it starts with the foundation and then progresses to the autoscaling of
worker nodes before moving on to example deployments of various applications in the
cluster. \\ %
The technology and programming languages utilized in the implementation are heterogeneous.
Nevertheless, according to a shared set of APIs and pre-defined structures, they
interact as a distinct and homogenous entity that maintains the entire cluster
in an active and healthy state. \\ %
Furthermore, several of the components and/or technologies discussed in this chapter
are interchangeable with other solutions. The latter is highly valuable for
organizations since it provides more configuration freedom and final control
over the cluster while continuously reflecting its ultimate goal.

\section{Distributions}
\label{sec:implementation_distributions}

% TODO Reference introduction dependencies

Almost all software relies on the necessity for a stable and consistent
operating system (OS). \\ %
The OS, as specified in section \ref{subsec:architecture_components_node}, is
installed on every node in the cluster and must be based on the \texttt{Linux
Kernel}. The latter is a critical requirement because most of the technologies
employed in the implementation rely on Linux primitives and core functionality
to operate properly. To enforce resource management for pods and containers, for
example, all Kubernetes distributions require the \texttt{cgroup v2}\footnote{\url{https://www.kernel.org/doc/html/latest/admin-guide/cgroup-v2.html}}
kernel module\footnote{\url{https://kubernetes.io/docs/concepts/architecture/cgroups}}.
Furthermore, the Linux Kernel is an Open Source project to which everyone can
contribute and get access, representing one fundamental objective of the entire project.
\\ %
To represent the potential of the implementation architecture being compatible with
several Linux distributions, the section is titled Distributions rather than
Operating System. A Linux distribution, often known as a Linux distro, is an operating
system that is built from components developed by various open-source projects and
programmers. The Linux kernel (the operating system's core), the GNU shell utilities
(terminal interface and commands), the X server (for a graphical desktop), the desktop
environment, a package management system, an installer, and other services are
all included in each distribution. Many components are created independently and
released in source code format. Thousands of software packages, utilities, and
applications can be found in a single Linux distribution. Linux distributions combine
code from open-source projects into a single operating system that can be installed
and booted. Linux distributions are available for a wide range of specialized
use cases, including desktop computers, servers without a graphical interface, supercomputers,
mobile devices, and more\cite{linux_distro}. \\ %
To be completely compatible with the implementation architecture, a Linux distribution
must meet three major requirements. They are described in the following list:
\begin{enumerate}
  \item The Linux Kernel version must be compatible with the Kubernetes version.
    \newline
    Newer versions of K8s may require newer versions of the Linux Kernel to
    provide newer functions and/or improve overall security.

  \item All necessary packages, as explained and documented in the section
    \ref{subsec:implementation_distributions_packages}, must be installed and globally
    accessible in the system.
    \newline
    A check on all essential packages is performed before the execution of the
    installation program, as detailed in section \ref{sec:implementation_installer}.
    If even one package is missing from the system, the installation fails with an
    error message identifying the missing packages.

  \item An Init System that is supported and compatible, as defined in the
    section \ref{subsec:implementation_distributions_init_system}.
    \newline
    OpenRC and Systemd are supported by the implementation. They are currently
    the two most used init systems, and the majority of available distributions are
    based on one of them.
\end{enumerate}
Finally, two custom Linux distributions are already configured, built, and
packaged as ISO files, as shown in section \ref{subsec:implementation_distributions_iso_image}.
They are based on a simple command line with no graphical interface and contain
only what is required to start a node. This is done to simplify the overall setup
by eliminating unnecessary software and technologies, addressing any
incompatibilities, and increasing the cluster's implementation and testing efficiency.
Both are derived from pre-existing Linux distributions that are known to be
stable, easy to configure, and lightweight, requiring just the bare minimum of resources
to operate. Furthermore, they are built on separate init systems, requiring the
overall implementation to be compatible with and tested with both.

\subsection{Packages}
\label{subsec:implementation_distributions_packages}

The packages used in the implementation that must be available in a Linux distribution
are specified in this section. A package is a collection of files that have been
bundled together and may be installed and uninstalled as a group\cite{package_manager}.
\\ %
All packages may be installed automatically using the package manager of the respective
Linux distribution. A package manager's job is to provide an interface that
assists the user in managing the collection of packages installed on the system.
A package manager maintains track of what software is installed on the system and
makes it simple to install new software, update software to newer versions, or delete
previously installed software. A package is usually a single piece of software. However,
it is typical for applications to be made up of numerous dependent packages. For
example, the \texttt{jq} package includes not just the \texttt{jq} command-line
JSON processor, but also the \texttt{oniguruma}\footnote{\url{https://github.com/kkos/oniguruma}}
package, a regular expressions library. Furthermore, various optional add-on packages
including extra data and documentation are available\cite{package_manager}. \\ %
All packages are already included by default in the two custom Linux
distributions. As a result, there is no need for any changes or the installation
of extra packages. \\ %
In the table \ref{tbl:packages}, each package has a brief overview taken from its
official documentation, as well as a description of where and why it has been
used. Furthermore, the required column indicates with symbols if the package is
required (\textcolor{bulmagreen}{\faicon{check}}), optional (\textcolor{bulmablue}{\faicon{question}}),
or not required (\textcolor{bulmared}{\faicon{times}}).

\begin{xltabular}
  {\textwidth} { >{\ttfamily}l | c | X | X }

  \multicolumn{1}{ c |}{\large{\textbf{Name}}} &
  \multicolumn{1}{ c |}{\large{\textbf{Required}}} &
  \multicolumn{1}{ c }{\large{\textbf{Overview}}} &
  \multicolumn{1}{ c }{\large{\textbf{Description}}} \\ \hline \hline

  coreutils & \textcolor{bulmagreen}{\faicon{check}} & & \\ \hline

  ethtool & \textcolor{bulmagreen}{\faicon{check}} & & \\ \hline

  inotify-tools & \textcolor{bulmagreen}{\faicon{check}} & & \\ \hline

  iproute2 & \textcolor{bulmagreen}{\faicon{check}} & & \\ \hline

  jq & \textcolor{bulmagreen}{\faicon{check}} & & \\ \hline

  ncurses & \textcolor{bulmared}{\faicon{times}} & & \\ \hline

  nodejs & \textcolor{bulmablue}{\faicon{question}} & & \\ \hline

  npm & \textcolor{bulmablue}{\faicon{question}} & & \\ \hline

  postgresql & \textcolor{bulmablue}{\faicon{question}} & & \\ \hline

  procps & \textcolor{bulmared}{\faicon{times}} & & \\ \hline

  sudo & \textcolor{bulmagreen}{\faicon{check}} & & \\ \hline

  sysbench & \textcolor{bulmagreen}{\faicon{check}} & & \\ \hline

  tzdata & \textcolor{bulmagreen}{\faicon{check}} & & \\ \hline

  util-linux & \textcolor{bulmagreen}{\faicon{check}} & & \\ \hline

  yq & \textcolor{bulmagreen}{\faicon{check}} & & \\ \hline

  \caption{Packages list}
  \label{tbl:packages}
\end{xltabular}

\subsection{Init System}
\label{subsec:implementation_distributions_init_system}

\subsubsection{OpenRC}
\label{subsec:implementation_distributions_init_system_openrc}

\subsubsection{Systemd}
\label{subsec:implementation_distributions_init_system_systemd}

\begin{wrapfigure}
  {r}{.25\textwidth}
  \centering
  \includegraphics[width=.25\textwidth]{images/logos/systemd.png}
  \caption{Systemd logo}
\end{wrapfigure}

\subsection{ISO image}
\label{subsec:implementation_distributions_iso_image}

\subsubsection{Generation}
\label{subsubsec:implementation_distributions_iso_generation}

\subsubsection{Alpine Linux}
\label{subsubsec:implementation_distributions_iso_alpine_linux}

\begin{wrapfigure}
  {r}{.25\textwidth}
  \centering
  \includegraphics[width=.25\textwidth]{images/logos/alpine.png}
  \caption{Alpine Linux logo}
\end{wrapfigure}

Alpine Linux\footnote{\url{https://www.alpinelinux.org}} is a lightweight, security-focused,
and resource-efficient distribution. In recent years, has grown in popularity as
the foundation for the majority of Docker images. It is based on musl\footnote{\url{https://musl.libc.org}}
(an implementation of the C standard library, an alternative to glibc), BusyBox\footnote{\url{https://busybox.net}}
(a single, small executable that combines tiny versions of many common UNIX utilities),
a custom package manager called APK\footnote{\url{https://wiki.alpinelinux.org/wiki/Alpine_Package_Keeper}}
and the OpenRC init system ()\cite{alpine_linux}. \\ % TODO OpenRC reference
Alpine Linux is the preferred reCluster distribution. It has been utilized and
tested during development and is now present in every cluster node.

\subsubsection{Arch Linux}
\label{subsubsec:implementation_distributions_iso_arch_linux}

\begin{wrapfigure}
  {r}{.25\textwidth}
  \centering
  \includegraphics[width=.25\textwidth]{images/logos/arch.png}
  \caption{Arch Linux logo}
\end{wrapfigure}

Arch Linux\footnote{\url{https://archlinux.org}} is a general-purpose
distribution that focuses on simplicity, minimalism, and code elegance. It is based
on Systemd init system () % TODO Systemd reference
and strives to stay bleeding edge, offering the latest stable versions of most
software through Pacman\footnote{\url{https://wiki.archlinux.org/title/Pacman}}
package manager. Uses a rolling release system that allows one-time installation
and perpetual software upgrades, allowing to not reinstall or upgrade the system
from one version to the next. Notably, Arch Linux is the foundation for several
many popular enterprise-grade distributions, such as Manjaro\footnote{\url{https://manjaro.org}}
and SteamOS\footnote{\url{https://store.steampowered.com/steamos}},
demonstrating its customizability, power, and stability\cite{arch_linux}.

\section{Dependencies}
\label{sec:implementation_dependencies}

% TODO Dependencies presentation

\subsection{Air-Gap Environment}
\label{subsec:implementation_dependencies_air_gap_environment}

\subsection{Management}
\label{subsec:implementation_dependencies_management}

\subsubsection{Configuration}
\label{subsec:implementation_dependencies_management_configuration}

\subsubsection{Script}
\label{subsec:implementation_dependencies_management_script}

\subsection{K3s}
\label{subsec:implementation_dependencies_k3s}

\begin{wrapfigure}
  {r}{.25\textwidth}
  \centering
  \includegraphics[width=.25\textwidth]{images/logos/k3s.png}
  \caption{K3s logo}
\end{wrapfigure}

\subsection{Node Exporter}
\label{subsec:implementation_dependencies_node_exporter}

\subsection{Prometheus}
\label{subsec:implementation_dependencies_prometheus}

\begin{wrapfigure}
  {r}{.25\textwidth}
  \centering
  \includegraphics[width=.25\textwidth]{images/logos/prometheus.png}
  \caption{Prometheus logo}
\end{wrapfigure}

\subsection{Autoscaler}
\label{subsec:implementation_dependencies_autoscaler}

\subsubsection{Cluster Autoscaler}
\label{subsec:implementation_dependencies_autoscaler_cluster_autoscaler}

\section{Installer}
\label{sec:implementation_installer}

\subsection{POSIX Shell}
\label{subsec:implementation_installer_posix_shell}

\subsection{Configuration Arguments}
\label{subsec:implementation_installer_configuration_arguments}

\subsection{Configuration Files}
\label{subsec:implementation_installer_configuration_files}

\subsubsection{Certificates}
\label{subsubsec:implementation_installer_configuration_files_certificates}

\subsubsection{K3s}
\label{subsubsec:implementation_installer_configuration_files_k3s}

\subsubsection{K8s}
\label{subsubsec:implementation_installer_configuration_filesn_k8s}

\subsubsection{Node exporter}
\label{subsubsec:implementation_installer_configuration_files_node_exporter}

\subsubsection{reCluster}
\label{subsubsec:implementation_installer_configuration_files_recluster}

\subsubsection{SSH}
\label{subsubsec:implementation_installer_configuration_files_ssh}

\subsection{Node Facts}
\label{subsec:implementation_installer_node_facts}

\subsection{Node Information}
\label{subsec:implementation_installer_node_information}

\subsection{Node Benchmarks}
\label{subsec:implementation_installer_node_benchmarks}

\subsubsection{SysBench}
\label{subsubsec:implementation_installer_node_benchmarks_sysbench}

\subsection{Node Power Consumption}
\label{subsec:implementation_installer_node_power_consumption}

\subsubsection{CloudFree Smart Plug}
\label{subsubsec:implementation_installer_node_power_consumption_cloudfree_smart_plug}

\subsection{Node Registration}
\label{subsec:implementation_installer_node_registration}

\subsubsection{K8s Node Label And Name}
\label{subsubsec:implementation_installer_node_registration_k8s_node_label_and_name}

\subsection{Cluster Initialization}
\label{subsec:implementation_installer_cluster_initialization}

\subsubsection{Kubeconfig}
\label{subsubsec:implementation_installer_cluster_initialization_kubeconfig}

\subsubsection{Database}
\label{subsubsec:implementation_installer_cluster_initialization_database}

\subsubsection{Server}
\label{subsubsec:implementation_installer_cluster_initialization_server}

\subsubsection{Admin And Autoscaler Users}
\label{subsubsec:implementation_installer_cluster_initialization_admin_and_autoscaler_users}

\subsubsection{K8s}
\label{subsubsec:implementation_installer_cluster_initialization_k8s}

\subsection{Services}
\label{subsec:implementation_installer_services}

\subsubsection{Start}
\label{subsubsec:implementation_installer_services_start}

\subsubsection{Stop}
\label{subsubsec:implementation_installer_services_stop}

\subsubsection{Boot}
\label{subsubsec:implementation_installer_services_boot}

\subsubsection{Shutdown}
\label{subsubsec:implementation_installer_services_shutdown}

\section{Server}
\label{sec:implementation_server}

\subsection{Database}
\label{subsec:implementation_server_database}

\subsubsection{PostgreSQL}
\label{subsubsec:implementation_server_database_postgresql}

\begin{wrapfigure}
  {r}{.25\textwidth}
  \centering
  \includegraphics[width=.25\textwidth]{images/logos/postgresql.png}
  \caption{PostgreSQL logo}
\end{wrapfigure}

\subsubsection{Prisma}
\label{subsubsec:implementation_server_database_prisma}

\begin{wrapfigure}
  {r}{.25\textwidth}
  \centering
  \includegraphics[width=.25\textwidth]{images/logos/prisma.png}
  \caption{Prisma logo}
\end{wrapfigure}

\subsubsection{Schema}
\label{subsubsec:implementation_server_database_schema}

\subsection{GraphQL}
\label{subsec:implementation_server_graphql}

\subsubsection{GraphQL vs REST}
\label{subsubsec:implementation_server_graphql_graphql_vs_rest}

\subsubsection{Schema}
\label{subsubsec:implementation_server_graphql_schema}

\subsection{Node Pool}
\label{subsec:implementation_server_node_pool}

\subsection{Node Registration}
\label{subsec:implementation_server_node_registration}

\subsubsection{JSON Web Token}
\label{subsubsec:implementation_server_node_registration_json_web_token}

\subsection{Upscaling}
\label{subsec:implementation_server_upscaling}

\subsubsection{Wake-on-LAN}
\label{subsubsec:implementation_server_scale_up_wake_on_lan}

\subsection{Downscaling}
\label{subsec:implementation_server_downscaling}

\subsubsection{SSH}
\label{subsubsec:implementation_server_scale_up_ssh}

\section{Autoscaling}
\label{sec:implementation_autoscaling}

\subsection{Vertical Pod Autoscaler}
\label{subsec:implementation_autoscaling_vertical_pod_autoscaler}

\subsection{Horizontal Pod Autoscaler}
\label{subsec:implementation_autoscaling_horizontal_pod_autoscaler}

\subsection{Cluster Autoscaler}
\label{subsec:implementation_autoscaling_cluster_autoscaler}

\subsubsection{Cloud Provider}
\label{subsubsec:implementation_autoscaling_cluster_autoscaler_cloud_provider}

\subsubsection{Configuration}
\label{subsubsec:implementation_autoscaling_cluster_autoscaler_configuration}

% TODO Cloud Configuration

\subsubsection{Upscaling}
\label{subsubsec:implementation_autoscaling_cluster_autoscaler_upscaling}

\subsubsection{Downscaling}
\label{subsubsec:implementation_autoscaling_cluster_autoscaler_downscaling}